% !TeX spellcheck = ru_RU
% !TEX root = vkr.tex

\section{Постановка задачи}
\label{sec:task}

Целью работы является развитие и подготовка к выпуску библиотеки pysatl-mpest для анализа смесей вероятностных распределений. Для её достижения поставлены следующие задачи:

\begin{enumerate}
    \item\label{task:core} завершить интеграцию библиотеки с вычислительным ядром pysatl-core и проверить работу основных сценариев (раздел~\ref{subsec:task1}; октябрь--ноябрь 2026~года);
    \item\label{task:architecture} спроектировать архитектуру оценщиков и инициализаторов. На основе обзора литературы, возможностей ядра и требований выбранных экспериментов определить состав поддерживаемых методов и распределений. Для каждого метода указать, будет ли он использовать аналитическое обновление параметров или общий оптимизатор (раздел~\ref{subsec:task2}; ноябрь--декабрь 2026~года);
    \item\label{task:estimators} реализовать выбранные итеративные и прямые алгоритмы оценки параметров и построения начальных приближений для согласованного набора распределений (раздел~\ref{subsec:task3}; декабрь 2026~года -- январь 2027~года);
    \item\label{task:preprocessing-design} спроектировать предварительный анализ выборки и определить, какие характеристики структуры смеси он будет оценивать (раздел~\ref{subsec:task4}; февраль 2027~года);
    \item\label{task:preprocessing} реализовать выбранные методы предварительного анализа и оценить качество полученных предположений и их влияние на оценку параметров (раздел~\ref{subsec:task5}; февраль--март 2027~года);
    \item\label{task:release} подготовить документацию, примеры и пакет к выпуску в согласованном объёме (раздел~\ref{subsec:task6}; в течение года, завершение к маю 2027~года);
    \item\label{task:experiments} воспроизвести эксперименты из научных работ, для которых доступны данные или описан способ их получения либо генерации. Сопоставить результаты с опубликованными и проанализировать расхождения (раздел~\ref{sec:experiment}; март--апрель 2027~года).
\end{enumerate}

Обзор литературы и проверка доступности экспериментальных данных начинаются в первом семестре и продолжаются при проектировании и реализации библиотеки.
